\chapter*{引言}
\markboth{引言}{引言}
\addcontentsline{toc}{chapter}{引言}

\textbf{第一印象。}\textbf{\trFukayaCategory}（Fukaya category）是关联到\trSymplectic\trManifold{}的一种代数结构。它编码了\trLagrange\trSubmanifold{}的信息、其中任意两个相交的方式，还有\trPseudoholomorphic\trDisk（或\trPolygon）的信息。打包进单个\trObject{}的数据量之大，使\trFukayaCategory{}既迷人又有些令人生畏，从\trSymplectic\trTopology{}的角度看尤其如此。在往下讲之前，我们想就这一观察再多说几句。这将是一段非正式的讨论，面向非专业读者，只求让人对局势的大概模样有个印象；其中包含许多并不严格的陈述，取材范围也比全书其余部分稍宽。

取定一个\trSymplectic\trManifold $M^{2n}$。粗略地说，\trFukayaCategory{}的\trObject{}就是\textbf{\trLagrange\trSubmanifold}（Lagrangian submanifold）$L^n \subset M$。这类\trSubmanifold{}构成的空间是无穷维的（局部上以$L$上的闭1-\trForm{}空间为模型）。然而，相互\trHamiltonianIsotopy{}的\trLagrange\trSubmanifold{}给出\trFukayaCategory{}中\trIsomorphic{}的\trObject，这把局部自由度削减到第一\trBettiNumber $b_1(L)$。的确，按\trHamiltonianIsotopy{}分类在若干简单情形中已被完全弄清：\trSurface{}上的曲线，那里问题本质上是一个\trTopology{}问题；不那么明显地，还有$S^2 \times S^2$中的\trLagrange\trSphere [75]，以及少数相关例子。另有一些（一般性的）效应对我们有利：并非所有\trLagrange\trSubmanifold{}都真正出现在\trFukayaCategory{}里——由于[60]意义下的\trObstruction；另一些可能出现却很平凡（即零\trObject）；而\trFukayaCategory{}中的\trIsomorphism{}也可能比\trLagrange\trIsotopy{}更弱。关于第一点，只需看看$S^2$中的\trCircle：所需的\trObstruction{}抵消只有在补集的两个半球的面积相等时才会发生；更有意思的情形见[31]、[32]。第二种现象出现在\trRealLocus{}$\RP^n \subset \CP^n$（$n>1$）上，这里假设它是Spin的（$n=4k-1$），并且\trFukayaCategory{}取$\QQ$-\trCoefficients。至于第三点，有一些形如$\CP^2$中\trLagrange\trTorus{}的候选例子，但就作者所知，还没有一个得到过严格验证。有了这些准备，对\trSymplectic\trTopology{}相对简单的\trManifold，原则上已能计算整个\trFukayaCategory；上文所引的论文（以及另一些，例如[132]、[102]、[101]、[59]）可以看作这个方向上更大规模的持续努力的一部分。尽管如此，对任何相当复杂的\trSymplectic\trManifold（例如$\CP^3$中\trDegree $\geq 4$ 的复\trHypersurface），直接分类\trFukayaCategory{}中\trObject{}的做法看来仍遥不可及。

说不定我们只是把目标设得太高了。毕竟，即便考虑纯代数的情境——比如有限维代数上的模——能实现完全分类的情形也只是极少数[41]。不过，这里有一个重要差别：对两个显式给出的模，有一个显而易见的算法可以判定它们是否\trIsomorphic。更一般地，给定有限个模，确定它们之间的\trMorphism{}以及\trMorphism{}的\trComposition{}映射都很直接。相比之下，在\trFukayaCategory{}中，\trMorphism{}是\trFloerCochainGroup，而\trComposition{}映射是这类群上的操作，两者都用适当的\trPseudoholomorphic{}映射\trModuliSpace{}来定义。有一些情形可以让这些空间得到几何的处理：当$M$是\trCotangent{}丛、而\trLagrange\trSubmanifold $L=\operatorname{graph}(df)$是\trSection{}时，一种“\trAdiabaticLimit”方法把问题化归为有限维\trMorseTheory [58]（这一方法的种种变体已在若干相关情形中取得成功）；当$M$是定义在实数上的\trAlgebraicVariety、而只考虑\trRealLocus{}$L=M_RR$时，可以应用代数几何工具；最后，当$M$是\trSurface{}时，对\trPseudoholomorphicCurve{}的研究在很大程度上变成组合性的。然而在这些情形之外，\trFloer{}理论的计算历来依赖逐个案例的办法，它们不提供任何\trApriori{}的成功保证。本书就是为处理这一问题而做的一次尝试。连同其他进展一道，它有望帮助把\trFukayaCategory{}变成一件更易于驾驭的工具。

\textbf{策略。}鉴于上文所述的困难，显然若想取得任何进展，就需要一些大幅度的简化。首先，我们只考虑\textbf{\trExactSymplectic\trManifold}（exact symplectic manifold）$M$，即\trSymplectic\trForm{}$\phi$是某个1-\trForm{}$\theta$的外\trDifferential{}的那些\trManifold。这当然意味着$M$不能是\trClosed{}的，但我们会施加标准的凸性条件，使这一缺陷对有限面积的\trPseudoholomorphicCurve{}不可见。\trAffineVariety{}是例子的一个好来源。$L \subset M$中最简单的一类\trLagrange\trSubmanifold{}是\trExact{}的那些，即$\theta|L$是某个函数的\trDifferential；我们只处理这些，其余一概\trApriori{}排除。这一限制有许多有用的技术后果：例如，\trExact\trLagrange\trSubmanifold{}不可能被\trObstruction{}拦住，并且其中每一个都是\trFukayaCategory{}的非零\trObject，其\trEndomorphism{}环是经典的\trCohomology{}环（没有量子修正）。除此之外，\trExact{}性更根本的好处在于它消除了\trFukayaCategory{}对\trSymplectic{}面积（或Novikov环）参数的依赖。\trMirrorSymmetry{}暗示这种依赖常常由非平凡的超越函数给出，一些示例计算（例如[108]、[57]）证实了这一点。在我们的设定里，这个参数被取为零（或无穷大，视理解而定），这使得情形从枚举的观点看乏味一些，却更便于计算。应当提到，归根结底，人们可以尝试把\trAffineVariety{}的\trFukayaCategory{}用作垫脚石，走向其\trProjective\trCompletion{}的\trFukayaCategory{}[129]。

为了更接近经典的\trHomology{}代数，我们还将施加一个Calabi--Yau型条件（确切地说，即$2c_1(M)=0$），并对\trLagrange\trSubmanifold{}施加相应的假设，这使得我们能给所有\trFloer{}群配上$\ZZ$-\trGrading。由此得到的框架可以总结如下。\trFukayaCategory{}$\scr{F}(M)$的\trObject——我们称之为\textbf{\trExactLagrangianBrane}（exact Lagrangian brane）——是三元组$L^\#=(L,\alpha^\#,P^\#)$，其中$L$是$M$中\trClosed{}的\trExact\trLagrange\trSubmanifold，$\alpha^\#$是$L$上一个实值函数，称为\trGrading（grading），而$P^\#$是$L$上的一个\trPinStructure。\trMorphism{}空间是\trFloerCochainGroup{}$CF^*(L_0^\#,L_1^\#)$，它们是有限维\trGraded\trVectorSpace（\trCoefficients{}域$\KK$取定但任意；\trPinStructure{}其实只在$\mathrm{char}(\KK)\neq 2$时才真正必要）。这些群上带有\trMultilinear\trComposition{}映射$\mu^d$（$d \geq 1$），它们构成一个\trAinfty{}结构。这个定义需要辅助的选择，但其结果在相差\trQuasiisomorphism{}的意义下与这些选择无关。

我们的第二个主要想法来自\trKontsevich{}的工作：取\trFukayaCategory{}的形式扩大。这些是纯代数的构造，可以应用于任何\textbf{\trACategory}（$A_\infty$-category）$\scr{A}$。第一步是引入\trChainComplex{}的类似物，即所谓的\trTwistedComplex（twisted complex），它们构成另一个包含$\scr{A}$的\trACategory{}$Tw\scr{A}$。它包含\trMorphism{}的\trMappingCone，因此正是我们所说的\trTriangulated\trACategory；由此得到的\trCohomology{}层面\trCategory{}$H^0(Tw\scr{A})$——也称作\trDerivedCategory{}$D(\scr{A})$——是经典意义下的\trTriangulated{}\trCategory。第二个有用的步骤是取\trDerivedCategory{}的\trSplitClosure（或\trKaroubi{}\trCompletion）$D^\pi(\scr{A})$，它为所有\trIdempotent\trEndomorphism{}引入形式\trDirectSummand；在\trAinfty{}层面上也有对应的构造，记作$\Pi(Tw\scr{A})$。有一个看似悖论、但其真确性对许多读者显而易见的陈述：每一次扩大都使这个\trCategory{}更容易描述。举一个例子——考虑\trSplitting\trGenerator{}的概念：若从一组\trObject{}出发，反复应用\trMappingCone{}并\trSplitting{}出\trDirectSummand，便能构造出全部\trObject，我们就说这组\trObject{}\trSplitGeneration{}了$\Pi(Tw\scr{A})$。每当$\scr{B} \subset \Pi(Tw\scr{A})$是由这样一组\trSplitting\trGenerator{}组成的\trFull{}的\trAinfty{}子\trCategory{}时，$\Pi(Tw\scr{B})$与$\Pi(Tw\scr{A})$都是\trQuasiequivalence{}的\trACategory。因此，若只关心$\Pi(Tw\scr{A})$，实际上确定一个这样的$\scr{B}$就够了。让我们记住，\trSplitGeneration{}是一条很弱的性质；比如说，它并不意味着\trGrothendieck{}群$K_0(D^\pi(\scr{A}))$由\trSplitting\trGenerator{}的类张成。因此，指望能找到有限多个\trLagrange\trSubmanifold{}来\trSplitGeneration{}$\Pi(Tw\scr{F}(M))$，这并非不合理；我们将会看到，在有利的情况下情形确实如此。

在与\trFukayaCategory{}相关联时，“计算”一词意味着底层的几何结构（\trSymplectic\trManifold、\trLagrange\trSubmanifold）必须以某种组合的方式被编码。为此，我们将系统地使用\textbf{\trPicardLefschetzTheory}（Picard--Lefschetz theory）。总的说来，这一进路在\trSymplectic{}几何中的重要性很大程度上源自\trDonaldson{}的工作[38]；这里我们只用该理论最初等的方面，即\trMonodromy{}映射的\trSymplectic{}本性，以及\trLefschetzThimble（Lefschetz thimble）与\trVanishingCycle{}的\trLagrange{}本性。为了利用这些性质，我们将（循着\trKontsevich{}的一个想法）引入\trLefschetzFibration{}的\trFukayaCategory，它们在某种意义上介于\trTotalSpace{}的平常\trFukayaCategory{}与\trSmooth\trFiber{}的\trFukayaCategory{}之间。随后它们将被用来搭建一台做计算的机器——在[126]中首次以猜想的方式提出——其原理是维数归纳。

\textbf{计算\trFukayaCategory。}为使表述简单，我们将把自己限于\trAlgebraicVariety{}的情形。具体地说，设$X \subset \CC^N$是一个$(n+1)$维\trAffineVariety，它\trSmooth，并且在无穷远处也\trSmooth。后一个条件是说，其闭包$\bar{X} \subset \CP^N$依然\trSmooth，并且与无穷远处\trHyperplane{}$\CP^{N-1} = \CP^N \setminus \CC^N$\trTransverse。我们给$X$配上Fubini--Study\trKaehler\trForm{}的限制，这使它成为一个\trExactSymplectic\trManifold。此外，我们将假设\trCanonical{}丛$\mathcal{K}_{\bar{X}}$对某个$d \in \ZZ$\trIsomorphic{}于$\scr{O}_{\bar{X}}(-d)$，这使得$X$成为Calabi--Yau的。这些性质也都为一个\trGeneric{}的仿射\trHyperplaneSection{}$X^1 = X \cap \CC^{N-1}$所继承。取$\scr{F}(X)$与$\scr{F}(X^1)$这两个\trFukayaCategory，\trCoefficients{}取在某个域$\KK$中。我们先讨论$X^1$，因为那边的结论要好一些：

\textbf{\trhTheorem{} A.}\ \emph{设$d \neq 2$、$n \geq 1$且$\mathrm{char}(\KK) \neq 2$，则$\Pi(Tw\scr{F}(X^1))$可由对$X$应用\trPicardLefschetzTheory{}所得到的组合数据计算出来。}

这个陈述需要稍作解释。首先，说$\Pi(Tw\scr{A})$对某个\trACategory{}$\scr{A}$是可计算的，意思如下：可以明确写出另一个\trACategory{}$\scr{B}$，它在强意义上有限（\trObject{}只有有限多个，\trMorphism{}构成有限维\trVectorSpace，且充分高阶$d \gg 0$的\trComposition{}映射$\mu^d_{\scr{B}}$全部为零），使得$\Pi(Tw\scr{B})$与$\Pi(Tw\scr{A})$是\trQuasiequivalence{}的\trACategory。特别是，这意味着$D^\pi(\scr{B}) \cong D^\pi(\scr{A})$（作为经典意义下的\trTriangulated{}\trCategory）。注意，对\trCoefficients{}域$\KK$的一个合理选择（比如有限域，或$\QQ$），$\scr{B}$确实只含有限多的信息。人们可以把这些信息输入一个计算机程序，它通过显式地构造\trTwistedComplex{}及其\trIdempotent\trEndomorphism，将生成$D^\pi(\scr{B})$中\trObject{}的\trIsomorphism{}类的完整（无穷）列表，连同\trMorphism、\trComposition{}映射与\textbf{\trExactTriangle}（exact triangle）一起。当然，这本身并不能解决关于$D^\pi(\scr{B})$的任何重大结构性问题。至于“由\trPicardLefschetzTheory{}产生的组合数据”的含义，我们把它精确的解释留给§19，但大致想法是：看从$X$到$\CC$与$\CC^2$的\trLinear{}映射，并用\trBraid{}单值来编码那些映射的分歧数据。同样的过程可以应用于$X^1$，并反复进行，直到维数降到零。已知定义$X$的显式方程，这些数据可以借助消去理论（elimination theory）算法式地提取出来（尽管除了最简单的例子，复杂度往往令人望而却步）。

我们还需稍微看一眼这些结果的局限。$n \geq 1$排除了最低的非平凡维数，也就是$X^1$是\trRiemann\trSurface{}的情形。这其实纯粹是为了方便：\trRiemann\trSurface{}的情形可以用类似（事实上还要简单些）的方法处理，只是需要一些特别的调整，那会使讨论复杂化。$\mathrm{char}(\KK) \neq 2$这个条件是我们证明的一个技术性产物——该证明使用了双重\trBranchedCover{}与$\ZZ/2$作用；很有可能存在别的方法能解除这一限制。另一个要求$d \neq 2$则多少更本质一些，我们把对它的讨论推迟到后面（见\trhRemark{} 19.8）。

如前所说，对$\scr{F}(X)$本身的结果要弱一些。这并不出人意料：比如说，目前还不知道有任何方法能判定$X$是否含有\trClosed{}的\trExact\trLagrange\trSubmanifold，从而也无法判定$\scr{F}(X)$（及其\trDerivedCategory）是否非平凡。这与$X^1$的情形恰成对照，在那里我们可以用\trVanishingCycle{}充当\trObject{}的一个天然来源。我们所能得到的最好结果是：

\textbf{\trhTheorem{} B.}\ \emph{假设$\mathrm{char}(\KK) \neq 2$。那么，从对$X$应用\trPicardLefschetzTheory{}所得的组合数据出发，可以构造出一个\trACategory，它包含一个\trQuasiequivalence{}于$\scr{F}(X)$的\trFull{}子\trCategory。}

这里所说的\trACategory{}形如$Tw\scr{B}$，其中$\scr{B}$在与之前相同的意义下有限。事实上，按构造$\scr{B}$将是\trDirected{}的，而这带来一些值得注意的结论。首先，$Tw\scr{B}$已经是\trSplitting{}闭的，因此它包含的将不只是$\scr{F}(X)$，而是整个$\Pi(Tw\scr{F}(X))$。其次，\trDirected{}的\trACategory{}\trQuasiisomorphism{}于\trMorphism{}空间为有限维的dg\trCategory。由于这一性质会被\trTwistedComplex{}的\trCategory{}继承，我们有：

\textbf{\trhCorollary{} C.}\ \emph{在\trhTheorem{} B的情形下，$\scr{F}(X)$\trQuasiisomorphism{}于一个具有有限维$\hom$空间的dg\trCategory。}

\trhTheorem{} B的实际用处，取决于我们对映射$\scr{F}(X) \to Tw\scr{B}$的几何意义理解得有多好。对$X$中最简单的一类\trLagrange\trSubmanifold——即能表示为\trHyperplaneSection{}的\trLefschetzPencil{}之\trMatchingCycle（matching cycle）的\trLagrange\trSphere——这些问题都容易处理，于是人们确实得到了一个计算\trFloerCohomology{}群的算法（严格说来，$n=1$的情形在此本应排除，理由与\trhTheorem{} A中相同）。对更一般的\trLagrange\trSubmanifold，这个陈述眼下还更多是理论性的。

上述结果的证明将在§19给出；它们会用到全书引入的大部分材料。作为一个自然的副产品，我们建立了\trVanishingCycle{}的\trHurwitz{}移动与\trDirected{}\trFukayaCategory{}的\trMutation{}之间的一般关系。事实上，我们会给出两个证明（\trhTheorem{} 17.20 与\trhTheorem{} 18.24；[127]中最初设想的证明与两者都不相同，目前仍未发表）。另一个有趣的结果是：对位于\trLefschetzFibration{}的\trTotalSpace{}中的\trLagrange\trSubmanifold，其\trFloerCohomology{}存在\trSpectralSequence（\trhCorollary{} 18.27）。当然，与任何\trSpectralSequence{}一样，这只不过把问题分解成一个可见的部分（起始项）与隐藏的高阶信息（\trDifferential）。在这个具体情形，\trDifferential{}表达的是\trLagrange\trSubmanifold{}如何由\trLefschetzThimble{}搭建起来，其中“搭建”一词按代数方式理解，即\trFukayaCategory{}中的Postnikov分解。

除结尾处的一处例外，本书不讨论例子与应用。幸运的是，现有文献弥补了这一欠缺（一些指引见第III部开头）。除此之外，任何对\trFukayaCategory{}真有严肃兴趣的读者，终究都得去啃开创性的[60]，它所讨论的许多几何与代数论题远远超出了本书建立的框架。
